\section{Rendering a Mathcraft Sheaf}\label{sec:rendering}

As we have established in \cref{sec:introduction}, Minecraft worlds with their
blocks and entities can be represented as sheaves $W \in \functorcat{\gamma^{\mathrm{op}}}{\Set}$
over a graph $\gamma$. The most natural next
question is: how do we \emph{render} such a sheaf?

% stick figures shared by the figures of this section
\tikzset{
	stick/.style={line width=0.5pt},
	% a person, standing upright
	pics/person/.style={code={
					\draw[stick] (-0.15cm,0cm) -- (0cm,0.4cm) -- (0.15cm,0cm)
					(0cm,0.4cm) -- (0cm,0.75cm)
					(-0.22cm,0.48cm) -- (0cm,0.68cm) -- (0.22cm,0.48cm);
					\draw[stick, fill=white] (0cm,0.87cm) circle[radius=0.12cm];
				}},
	% a zombie, arms stretched out in front of it
	pics/zombie/.style={code={
					\draw[stick] (-0.15cm,0cm) -- (0cm,0.4cm) -- (0.15cm,0cm)
					(0cm,0.4cm) -- (0cm,0.75cm)
					(0cm,0.69cm) -- (0.3cm,0.69cm)
					(0cm,0.62cm) -- (0.3cm,0.62cm);
					\draw[stick, fill=white] (0cm,0.87cm) circle[radius=0.12cm];
				}},
}

% paragraph on the left, Figure (entity sheaf) on the right
{%
	\tikzset{
		twochunks/.style={
				x={(1.0cm,0cm)}, y={(0.55cm,0.38cm)}, z={(0cm,1.0cm)},
				line join=round, line cap=round},
		chunkline/.style={line width=0.4pt, dash pattern=on 2pt off 1.5pt},
		hiddenline/.style={line width=0.3pt, dash pattern=on 0.8pt off 1.6pt},
	}
	% the list of entities at a point: brackets of half-width #2 (in cm)
	% placed just above the coordinate #1, with contents #3
	\newcommand{\entitylist}[3]{%
		\path (#1) ++(0cm,0.12cm) coordinate (#1-label);
		\begin{scope}[shift={(#1-label)}, scale=1/3, every node/.style={scale=1/3}]
			\fill[white] (-#2cm,-0.1cm) rectangle (#2cm,1.1cm);
			\draw[line width=0.45pt]
			({-#2cm+0.08cm},-0.08cm) -- (-#2cm,-0.08cm) -- (-#2cm,1.08cm) -- ({-#2cm+0.08cm},1.08cm)
			({#2cm-0.08cm},-0.08cm) -- (#2cm,-0.08cm) -- (#2cm,1.08cm) -- ({#2cm-0.08cm},1.08cm);
			#3
		\end{scope}
	}
	\noindent
	\begin{minipage}[t]{0.58\textwidth}
		\vspace{0pt}
		Rendering here means that we want to draw an image onto a (computer) screen
		that displays the information contained in $W$. During gameplay, a user is
		typically assigned a dedicated entity called the \emph{player}, and the user
		has a \emph{rendering distance}. The player is given by a unique
		$e \in \EntityID$ contained in the image of the entity sheaf $E$ of $W$, i.e.\
		$e \in \operatorname{flatten} \im E(\gamma)$. We can find at least one chunk
		$C_{i,j,k}$ such that $e \in E(C_{i,j,k})$, and then $e_p = (x,y,z)$, where
		$(x,y,z) \in [0,w) \times [0,d) \times [0,h)$ is uniquely determined by
		$e \in E(C_{i,j,k})(x,y,z)$. By its $\EntityID$, we wish to associate
		with it a velocity $e_v = \partial e_p / \partial t$, an acceleration
		$e_a = \partial^2 e_p / \partial t^2$, a \emph{heading
			direction} $e_d \in S^2$, which encodes in what direction the entity looks,
		and an \emph{entity height} $e_h \in \RR_{\ge 0}$, which is the offset
		between $e_p$ and where the camera shall be.
	\end{minipage}%
	\hfill
	\begin{minipage}[t]{0.4\textwidth}
		\vspace{0pt}
		\centering
		\begin{tikzpicture}[twochunks]
			\def\w{2.2}\def\d{1.6}\def\h{1.8}
			% hidden edges
			\draw[hiddenline] (0,\d,0) -- (2*\w,\d,0)
			(0,\d,0) -- (0,0,0)
			(0,\d,0) -- (0,\d,2*\h)
			(\w,0,0) -- (\w,\d,0) -- (\w,\d,2*\h)
			(0,0,\h) -- (0,\d,\h) -- (2*\w,\d,\h)
			(\w,0,\h) -- (\w,\d,\h);
			% visible edges
			\draw[chunkline] (0,0,0) -- (2*\w,0,0) -- (2*\w,0,2*\h) -- (0,0,2*\h) -- cycle
			(\w,0,0) -- (\w,0,2*\h) -- (\w,\d,2*\h)
			(0,0,\h) -- (2*\w,0,\h) -- (2*\w,\d,\h)
			(0,0,2*\h) -- (0,\d,2*\h) -- (2*\w,\d,2*\h) -- (2*\w,0,2*\h)
			(2*\w,0,0) -- (2*\w,\d,0) -- (2*\w,\d,2*\h);
			% chunk names
			\node[below] at (0.5*\w,0,0) {$C_{i-w,j,k}$};
			\node[below] at (1.5*\w,0,0) {$C_{i,j,k}$};
			\node[fill=white, inner sep=1pt] at (0.5*\w,0,1.5*\h) {$C_{i-w,j,k+h}$};
			\node[fill=white, inner sep=1pt] at (1.5*\w,0,1.5*\h) {$C_{i,j,k+h}$};
			% the player e at e_p = (x,y,z), measured from the corner of C_{i,j,k}
			\def\px{0.7}\def\py{0.8}\def\pz{0.6}
			\coordinate (player) at (\w+\px,\py,\pz);
			\begin{scope}[line width=0.5pt, -{Stealth[length=4pt]}]
				\draw (\w,0,0) -- node[above, inner sep=1.5pt] {$x$} (\w+\px,0,0);
				\draw (\w+\px,0,0) -- node[above left, inner sep=0.5pt] {$y$} (\w+\px,\py,0);
				\draw (\w+\px,\py,0) -- node[right, inner sep=1.5pt] {$z$} (player);
			\end{scope}
			% two zombies at another point of C_{i,j,k}, one zombie in C_{i-w,j,k}
			\coordinate (zombies) at (\w+1.5,1.1,0.7);
			\coordinate (zombie) at (0.8,0.9,0.7);
			\draw[hiddenline] (zombies) -- (\w+1.5,1.1,0) (zombie) -- (0.8,0.9,0);
			\fill (player) circle[radius=1.4pt];
			\fill (zombies) circle[radius=1.4pt];
			\fill (zombie) circle[radius=1.4pt];
			\entitylist{player}{0.35}{\pic at (0cm,0cm) {person};}
			\entitylist{zombies}{0.72}{%
				\pic at (-0.38cm,0cm) {zombie};
				\node[inner sep=0pt, scale=2] at (0.08cm,0.1cm) {,};
				\pic at (0.3cm,0cm) {zombie};
			}
			\entitylist{zombie}{0.4}{\pic at (-0.07cm,0cm) {zombie};}
		\end{tikzpicture}
		\captionof{figure}{Depiction of the entity sheaf and how coordinates are calculated.}
		\label{fig:entity-lists}
	\end{minipage}
	\par\medskip
}

The rendering distance in classical Minecraft is given by some natural number
$r \in \NN$. Given the chunk $C_i$ the player is in, it considers the open set
\[
	U = \bigcup_{\operatorname{dist}(i,j) \le r} C_j,
\]
% paragraph on the left, top-down view of the rendered chunks on the right
{%
	\tikzset{
		topview/.style={x=0.165cm, y=0.165cm, line join=round},
		blk/.style={line width=0.2pt},
		grass/.style={pattern={Dots[distance=2.2pt, radius=0.28pt]}},
		water/.style={pattern={Lines[angle=0, distance=1.1pt, line width=0.3pt]}},
		leaves/.style={pattern={Hatch[angle=0, distance=1.1pt, line width=0.3pt]}},
		stone/.style={pattern={Dots[distance=1.3pt, radius=0.3pt]}},
		wood/.style={fill=black},
	}
	% a block at (#1,#2) in block coordinates, styled #3
	\newcommand{\topblock}[3]{\draw[blk, #3] (#1,#2) rectangle ++(1,1);}
	% a tree: 3x3 leaves around the trunk at (#1,#2)
	\newcommand{\toptree}[2]{%
		\foreach \u in {-1,0,1} {\foreach \v in {-1,0,1} {\topblock{#1+\u}{#2+\v}{leaves}}}
		\topblock{#1}{#2}{wood}%
	}
	% an entity #3 standing at (#1,#2), on a small white background
	\newcommand{\topentity}[3]{%
		\fill[white] (#1-0.9,#2-0.2) rectangle (#1+1.1,#2+1.75);
		\pic[scale=1/3] at (#1,#2) {#3};
	}
	\par\vspace{-\parskip}\noindent
	\begin{minipage}[t]{0.58\textwidth}
		\vspace{0pt}
		where $\operatorname{dist} \colon \gamma \times \gamma \to \NN$ counts how far
		apart two nodes of the graph are, that is, the length of a shortest path
		between them, disregarding any nodes that do not correspond to chunks, such as
		those for faces. Thus two chunks are one step apart exactly when they share a
		common face. For the lattice $X$ from \cref{fig:chunk-lattice},
		$\operatorname{dist}$ is then just the taxicab metric on the chunk indices
		$(i/w, j/d, k/h) \in \ZZ^3$.

		\medskip
		We may however choose any open set, or subgraph, $U$ for rendering, and we
		will call it the \emph{rendering set}.
	\end{minipage}%
	\hfill
	\begin{minipage}[t]{0.4\textwidth}
		\vspace{0pt}
		\centering
		\begin{tikzpicture}[topview]
			% the world from above: 9 x 5 chunks of 4 x 4 blocks each,
			% the player standing in the centre chunk (4,2)
			\fill[grass] (0,0) rectangle (36,20);
			% a lake
			\foreach \x/\y in {25/3,26/3,27/3,24/4,25/4,26/4,27/4,28/4,24/5,25/5,26/5,27/5,28/5,29/5,25/6,26/6,27/6,28/6,26/7}
				{\topblock{\x}{\y}{water}}
			% a stone path
			\foreach \x in {0,...,13} {\topblock{\x}{10}{stone}}
			\foreach \y in {11,...,14} {\topblock{13}{\y}{stone}}
			% trees
			\toptree{3}{14}
			\toptree{10}{5}
			\toptree{21}{16}
			\toptree{32}{12}
			\toptree{33}{4}
			% chunk borders, dotted, except between two chunks of the rendering set
			\begin{scope}[line width=0.3pt, dash pattern=on 0.8pt off 1.2pt]
				\foreach \a in {0,...,9} {\foreach \b in {0,...,4} {
							\pgfmathtruncatemacro{\keep}{
								(abs(\a-5)+abs(\b-2) > 2) || (abs(\a-4)+abs(\b-2) > 2)}
							\ifnum\keep=1 \draw (4*\a,4*\b) -- ++(0,4); \fi
						}}
				\foreach \a in {0,...,8} {\foreach \b in {0,...,5} {
							\pgfmathtruncatemacro{\keep}{
								(abs(\a-4)+abs(\b-3) > 2) || (abs(\a-4)+abs(\b-2) > 2)}
							\ifnum\keep=1 \draw (4*\a,4*\b) -- ++(4,0); \fi
						}}
			\end{scope}
			% the rendering set U: chunks within taxicab distance 2
			\def\rendersetpath{(16,20) -- (20,20) -- (20,16) -- (24,16) -- (24,12)
				-- (28,12) -- (28,8) -- (24,8) -- (24,4) -- (20,4) -- (20,0)
				-- (16,0) -- (16,4) -- (12,4) -- (12,8) -- (8,8) -- (8,12)
				-- (12,12) -- (12,16) -- (16,16) -- cycle}
			% dither everything outside of it
			\fill[even odd rule, pattern={Lines[angle=45, distance=2pt, line width=0.35pt]}]
			(0,0) rectangle (36,20) \rendersetpath;
			\draw[line width=0.9pt] \rendersetpath;
			% the player in the centre chunk, and some zombies
			\topentity{18}{9.8}{person}
			\topentity{14.5}{12.5}{zombie}
			\topentity{22}{9.5}{zombie}
			\topentity{18}{13}{zombie}
			\topentity{2}{17.5}{zombie}
			\topentity{34}{1}{zombie}
			\topentity{29.5}{14}{zombie}
			\topentity{6}{1.5}{zombie}
		\end{tikzpicture}
		\captionof{figure}{The chunks within taxicab distance $2$ of the
			player (bold) form the rendering set $U$ for $r = 2$.}
		\label{fig:render-distance}
	\end{minipage}
	\par\medskip
}

To model rendering, we will use the topos of trees
$\toposoftrees \coloneqq \functorcat{\omega^{\mathrm{op}}}{\Set}$, which we
have introduced and studied in \cite{happelToposTrees2026}. In particular,
this topos allows us to model ray tracing in its most natural way, with the
added benefit that each render pass will be productive.
\todo[inline]{Possibly also add the topos for synthetic differential geometry here?}
The general idea is that we construct a term $\rho$ of the coinductive stream
type
\[
	\Stream(U) \cong U \times \later \Stream(U)
\]
in $\toposoftrees$, which starts at some point $p$ in a chunk $C_j \subseteq U$
by setting
\[
	\restr{\rho}{0} = [\rho_0] \coloneqq [p],
\]
and which follows a geodesic spanned by a tangent vector $v_0 \in T_p C_j$
until it collides with a block or an entity. Here
$\exp \colon (T_p C_j, +) \to (\RR^3, +)$ is the exponential map of Lie groups
\cite{happelExponentialMapLie2026}. Using $\exp$, we can define
\[
	\lambda_{n+1} \coloneqq \inf \{\lambda > 0 \mid \rho_n + \exp(\lambda v_n)
	\text{ lies in a block or an entity}\} \in \RR_{\ge 0} \cup \{\infty\}
\]
as the distance (in units of $v_n$) to the next collision, so that
$\rho_{n+1} = \rho_n + \exp(\lambda_{n+1} v_n)$ is the point of collision.
(The addition is well defined, as $U$ is covered by chunks $C_j$, which are each
diffeomorphic to $[0,w] \times [0,d] \times [0,h] \subseteq \RR^3$, and as such
can be identified with subsets of $\RR^3$.) There, the ray is reflected and continues in the direction
\[
	v_{n+1} \coloneqq v_n - 2 \langle v_n, \nu_{n+1} \rangle \nu_{n+1},
\]
where $\nu_{n+1}$ is the unit normal of the surface that was hit, which follows
from the law of reflection with respect to the tangent plane
$t = \nu_{n+1}^{\perp}$.

For blocks, colliding is easy, as the tangent plane $t$ is clear. Every face
of a block is axis-aligned, so $\nu_{n+1} \in \{\pm e_1, \pm e_2, \pm e_3\}$, and
the reflection simply flips the sign of one coordinate of $v_n$. Thus, when
colliding with a block,
\[
	\restr{\rho}{n+1} = [\rho_0, \dots, \rho_n, \rho_n + \exp(\lambda_{n+1} v_n)],
	\qquad
	v_{n+1} = v_n - 2 \langle v_n, \nu_{n+1} \rangle \nu_{n+1}.
\]
For entities, we would indeed need to know some additional information, such as
the shape of the entity. For simplicity, we will just assume that every entity
$f$ is a sphere of radius $s$ around its position $f_p$. The unit normal at the
point of collision is then $\nu_{n+1} = (\rho_{n+1} - f_p) / s$, which gives,
when colliding with an entity $f$,
\[
	\restr{\rho}{n+1} = [\rho_0, \dots, \rho_n, \rho_n + \exp(\lambda_{n+1} v_n)],
	\qquad
	v_{n+1} = v_n - \frac{2}{s^2} \langle v_n, \rho_{n+1} - f_p \rangle (\rho_{n+1} - f_p).
\]

Finally, assume that the ray starting at $\rho_n$ in direction $v_n$ never
collides with a block or an entity, or that it reaches the boundary
$\partial U$ of the rendering set before doing so. In both cases we set
$\lambda_{n+1} = \infty$, and the ray ends there, so we let the stream stay
constant from then on,
\[
	\restr{\rho}{n+1} = [\rho_0, \dots, \rho_n, \rho_n],
	\qquad
	v_{n+1} = v_n.
\]

We may therefore generalize this procedure of ray tracing to a family of
global elements in the topos of trees. Given a starting position $x \in U$ and
a heading vector $v \in T_x U$, this yields
\begin{align*}
	\mathrm{ray}_{x,v}         & \colon 1 \to \Stream(U), \\
	\mathrm{ray}_{x,v}         & \coloneqq [\rho_0, \rho_1, \dots],
\end{align*}
where $\rho$ is constructed as above, starting with $\rho_0 = x$ and
$v_0 = v$.

We can now consider the oriented plane of width $\mathrm{plane}_w$ and height
$\mathrm{plane}_h$, spanned by orthogonal unit vectors $u, v \in T_x U$,
given by
\[
	\Pi \coloneqq \left\{ x + \exp(a u + b v) \in U \;\middle|\;
	a \in \left[-\tfrac{\mathrm{plane}_w}{2}, \tfrac{\mathrm{plane}_w}{2}\right],\
	b \in \left[-\tfrac{\mathrm{plane}_h}{2}, \tfrac{\mathrm{plane}_h}{2}\right]
	\right\},
\]
where $\exp \colon (T_x U, +) \to (\RR^3, +)$ is the exponential map from
above.
We can then set $x = e_p + \delta$ for some offset $\delta$. By choosing another
point $y = e_p + \delta'$, we can then consider all rays from $y$ through the
points of $\Pi$. This gives us a natural choice of heading vector for
$\mathrm{ray}_{e_p + \delta', v}$: for every point $q \in \Pi$, we take
$v = q - y$, the direction from $y$ towards $q$.

\begin{figure}[ht]
	\centering
	% chunk grid and projection as in Fig.~\ref{fig:entity-sheaf}
	\tikzset{
		sheafpic/.style={
				x={(1.0cm,0cm)}, y={(0.55cm,0.38cm)}, z={(0cm,1.0cm)},
				scale=1.05, line join=round, line cap=round},
		chunkline/.style={line width=0.4pt, dash pattern=on 2pt off 1.5pt},
		chunkstub/.style={line width=0.3pt, dash pattern=on 0.8pt off 1.6pt},
		vox/.style={line width=0.3pt, preaction={fill=white}},
		% ordered "dithering": black and white only
		dither sparse/.style={pattern={Dots[distance=2.2pt, radius=0.28pt]}},
		dither mid/.style={pattern={Dots[distance=1.3pt, radius=0.3pt]}},
		dither lines/.style={pattern={Lines[angle=45, distance=1.3pt, line width=0.3pt]}},
		dither vert/.style={pattern={Lines[angle=90, distance=1.1pt, line width=0.3pt]}},
		dither dark/.style={pattern={Hatch[angle=45, distance=1.1pt, line width=0.3pt]}},
	}
	% a unit cube at voxel (#1,#2,#3), faces styled #4 (top), #5 (front), #6 (right)
	\newcommand{\vox}[6]{%
		\draw[vox,#5] (#1,#2,#3) -- ++(1,0,0) -- ++(0,0,1) -- ++(-1,0,0) -- cycle;
		\draw[vox,#6] ({#1+1},#2,#3) -- ++(0,1,0) -- ++(0,0,1) -- ++(0,-1,0) -- cycle;
		\draw[vox,#4] (#1,#2,{#3+1}) -- ++(1,0,0) -- ++(0,1,0) -- ++(-1,0,0) -- cycle;}
	\newcommand{\dirt}[3]{\vox{#1}{#2}{#3}{dither sparse}{dither mid}{dither dark}}
	\newcommand{\grass}[3]{\vox{#1}{#2}{#3}{dither sparse}{dither mid}{dither dark}}
	\newcommand{\wood}[3]{\vox{#1}{#2}{#3}{dither mid}{dither vert}{dither dark}}
	\newcommand{\leaves}[3]{\vox{#1}{#2}{#3}{dither sparse}{dither lines}{dither dark}}
	\newcommand{\chest}[3]{%
		\begin{scope}[shift={({#1+0.1},{#2+0.1},#3)}, scale=0.8]
			\vox{0}{0}{0}{}{}{dither sparse}
			\draw[line width=0.3pt] (0,0,0.62) -- (1,0,0.62) -- (1,1,0.62);
			\fill (0.42,0,0.48) rectangle (0.58,0,0.72);
		\end{scope}}
	% the chunk grid of Fig.~\ref{fig:chunk-tiling}, outlined by dashed lines only
	\newcommand{\chunkgrid}{%
		\def\w{1}\def\d{1}\def\h{1.25}
		\def\nx{4}\def\ny{3}\def\nz{2}
		\pgfmathsetmacro{\X}{\nx*\w}
		\pgfmathsetmacro{\Y}{\ny*\d}
		\pgfmathsetmacro{\Z}{\nz*\h}
		\foreach \i in {0,...,\nx} {
				\draw[chunkline] (\i*\w,0,0) -- (\i*\w,0,\Z) -- (\i*\w,\Y,\Z);
				\draw[chunkstub] (\i*\w,0,0) -- (\i*\w,0,-0.45);
				\draw[chunkstub] (\i*\w,\Y,\Z) -- (\i*\w,\Y+0.6,\Z);
			}
		\foreach \k in {0,...,\nz} {
				\draw[chunkline] (0,0,\k*\h) -- (\X,0,\k*\h);
				\draw[chunkstub] (0,0,\k*\h) -- (-0.45,0,\k*\h);
				\draw[chunkstub] (\X,\Y,\k*\h) -- (\X,\Y+0.6,\k*\h);
			}
		\foreach \j in {0,...,\ny} {
				\draw[chunkline] (0,\j*\d,\Z) -- (\X,\j*\d,\Z) -- (\X,\j*\d,0);
				\draw[chunkstub] (0,\j*\d,\Z) -- (-0.45,\j*\d,\Z);
				\draw[chunkstub] (\X,\j*\d,0) -- (\X,\j*\d,-0.45);
			}
		\draw[chunkline] (\X,0,0) -- (\X,\Y,0) (\X,0,\h) -- (\X,\Y,\h);
		% hidden edges of the distinguished chunk C_{w,0,h}
		\draw[chunkline] (\w,\d,\h) -- (2*\w,\d,\h);
		\draw[chunkline] (\w,0,\h) -- (\w,\d,\h) -- (\w,\d,2*\h);
		\draw[chunkline] (2*\w,0,\h) -- (2*\w,\d,\h) -- (2*\w,\d,2*\h);
		% dimensions of a single chunk
		\draw[{Bar[width=4pt]}-{Bar[width=4pt]}, line width=0.4pt]
		(0,0,-0.75) -- node[fill=white, inner sep=1.5pt] {$w$} (\w,0,-0.75);
		\draw[{Bar[width=4pt]}-{Bar[width=4pt]}, line width=0.4pt]
		(\X+0.55,0,-0.1) -- node[below right=-1pt, inner sep=1.5pt] {$d$} (\X+0.55,\d,-0.1);
		\draw[{Bar[width=4pt]}-{Bar[width=4pt]}, line width=0.4pt]
		(-0.85,0,0) -- node[fill=white, inner sep=1.5pt] {$h$} (-0.85,0,\h);
	}
	% voxel coordinates inside C_{w,0,h}: 5 x 5 x 6 cubes of side 0.2
	\tikzset{inchunk/.style={shift={(1.1,0,1.25+0.025)}, scale=0.2}}
	% the player at e_p = (2.5, 2.5, 2) in voxel coordinates of C_{w,0,h}
	% (5 x 5 x 6 cubes of side 0.2), and the plane Pi through e_p
	\newcommand{\playerwithplane}[2]{%
		\coordinate (ep) at (2.5,2.5,2);
		\fill[pattern={Dots[distance=#1, radius=0.3pt]}]
		(2.5,1,0.5) -- (2.5,4,0.5) -- (2.5,4,3.5) -- (2.5,1,3.5) -- cycle;
		\draw[line width=0.5pt]
		(2.5,1,0.5) -- (2.5,4,0.5) -- (2.5,4,3.5) -- (2.5,1,3.5) -- cycle;
		#2
		\begin{scope}[shift={(ep)}, line width=0.5pt]
			\draw (-0.3cm,0cm) -- (0cm,0.8cm) -- (0.3cm,0cm);
			\draw (0cm,0.8cm) -- (0cm,1.45cm);
			\draw (-0.45cm,0.9cm) -- (0cm,1.3cm) -- (0.45cm,0.9cm);
			\draw[preaction={fill=white}] (0cm,1.67cm) circle[radius=0.22cm];
		\end{scope}}
	\begin{subfigure}[b]{0.58\textwidth}
		\centering
		\begin{tikzpicture}[sheafpic]
			\chunkgrid
			\begin{scope}[inchunk]
				\playerwithplane{1.2pt}{}
			\end{scope}
		\end{tikzpicture}
		\caption{The player in the chunk $C_{w,0,h}$.}
		\label{fig:render-plane-overview}
	\end{subfigure}%
	\hfill
	\begin{subfigure}[b]{0.4\textwidth}
		\centering
		\begin{tikzpicture}[sheafpic, scale=3.2]
			% the chunk C_{w,0,h}, outlined by dashed lines
			\draw[chunkline] (0,0,0) -- (1,0,0) -- (1,0,1.25) -- (0,0,1.25) -- cycle;
			\draw[chunkline] (1,0,0) -- (1,1,0) -- (1,1,1.25) -- (1,0,1.25);
			\draw[chunkline] (0,0,1.25) -- (0,1,1.25) -- (1,1,1.25);
			\draw[chunkline] (0,0,0) -- (0,1,0) -- (1,1,0);
			\draw[chunkline] (0,1,0) -- (0,1,1.25);
			\node[font=\small, anchor=north west, inner sep=2pt] at (0,0,1.25) {$C_{w,0,h}$};
			\begin{scope}[shift={(0.1,0,0.025)}, scale=0.2]
				\playerwithplane{1.6pt}{%
					% tangent vectors at e_p, drawn behind the player
					\draw[-{Stealth[length=4pt, width=3pt]}, line width=0.6pt] (ep) -- ++(2,0,0)
					node[font=\small, anchor=west, inner sep=1.5pt] {$e_d$};
					\draw[-{Stealth[length=4pt, width=3pt]}, line width=0.6pt] (ep) -- ++(0,1.5,0)
					node[font=\small, anchor=south west, inner sep=1pt] {$u$};
					\draw[-{Stealth[length=4pt, width=3pt]}, line width=0.6pt] (ep) -- ++(0,0,2.6)
					node[font=\small, anchor=south, inner sep=1.5pt] {$v$};}
				\node[font=\small, anchor=south east, inner sep=1pt] at (2.5,1,3.5) {$\Pi$};
				\fill (ep) circle[radius=0.07];
				\node[font=\small, anchor=north west, fill=white, inner sep=1pt, xshift=2pt, yshift=-2pt] at (ep) {$x = e_p$};
			\end{scope}
		\end{tikzpicture}
		\caption{Zoomed into $C_{w,0,h}$.}
		\label{fig:render-plane-zoom}
	\end{subfigure}
	\caption{The plane $\Pi$ at the point $x$, here with $\delta = 0$, so that $x = e_p$ is the position of the player. It is spanned by the
		orthogonal unit vectors $u$ and $v$, here chosen orthogonal to the heading
		direction $e_d$.}
	\label{fig:render-plane}
\end{figure}

% TODO
